\section{Experimental Setup}

We design experiments to answer four research questions mapping to our assumptions:

\textbf{RQ1:} Can design features be extracted objectively at scale? (Tests A2: feature extraction protocol)

\textbf{RQ2:} Do coverage families capture meaningful groupings? (Tests A4: family similarity)

\textbf{RQ3:} Do families predict future adoption? (Tests main claim + A3: temporal persistence)

\textbf{RQ4:} Can NLP classify citation contexts reliably? (Tests A1: citation evidence quality)

\subsection{Datasets}

\textbf{Benchmark Paper Corpus:} 20 diverse benchmarks stratified across modalities:
\begin{itemize}
\item \textbf{Vision (5):} ImageNet~\cite{imagenet}, COCO~\cite{coco}, PASCAL VOC, CelebA, CIFAR-10
\item \textbf{Language-Translation (3):} WMT14, WMT16, IWSLT
\item \textbf{Language-QA (4):} SQuAD~\cite{squad}, Natural Questions, TriviaQA, HotpotQA
\item \textbf{Audio/Multimodal (8):} LibriSpeech, Common Voice, VGGSound, VoxCeleb, MS-COCO (multimodal), Conceptual Captions, Visual Genome, MSCOCO Captions
\end{itemize}

\textbf{Rationale:} Pilot sample (20 vs 100+ targeted corpus) enables proof-of-concept while covering major modalities. Stratified sampling ensures diversity for cluster discovery. Benchmarks selected with $\geq 50$ citations to ensure sufficient usage data.

\textbf{Citation Data:} Extracted via Semantic Scholar API for 2015--2024 papers citing selected benchmarks. Historical split: pre-2023 for training, 2023--2024 for test (temporal validation).

\textbf{Synthetic Citation Contexts (H-E1):} 100 template-generated validation claims vs baseline mentions for SciBERT classifier training. Real-world validation deferred to future work.

\subsection{Baselines}

\textbf{Random Baseline:} Random assignment of benchmarks to coverage families. Expected citation overlap $\leq 50\%$ (null hypothesis H0).

\textbf{Citation-Count Baseline:} Rank by popularity (total citations). Tests whether popularity alone predicts coverage.

\textbf{Manual Expert Review:} Gold standard for accuracy but requires weeks per hypothesis (not scalable). Comparison deferred to Phase 5.

\subsection{Evaluation Metrics}

\textbf{H-M1 (Feature Extraction):} Cohen's kappa for categorical features (task, modality, metrics), ICC for continuous (dataset size). Target: kappa $>0.80$.

\textbf{H-M2 (Clustering Quality):} Intra-family cosine similarity, silhouette score, modularity on citation network. Target: similarity $\geq 0.60$.

\textbf{H-M3 (Prediction Accuracy):} Jaccard similarity of citation sets within families. Target: $>70\%$ with $p<0.05$ vs random.

\textbf{H-E1 (Classification):} Precision/recall on held-out synthetic test set. Target: precision $>85\%$.

\subsection{Implementation Details}

\begin{itemize}
\item \textbf{Embeddings:} SentenceBERT (all-MiniLM-L6-v2) for benchmark descriptions
\item \textbf{Clustering:} K-means with $k \in \{2,4,6,8\}$, selected via silhouette score
\item \textbf{Citation Classifier:} SciBERT fine-tuned on synthetic data (80/20 train/test split)
\item \textbf{Statistical Testing:} Permutation test (1000 iterations) for citation overlap significance
\end{itemize}
